% META: {"id": "1SPE-SUITES-ME-005", "chapitre": "1SPE-SUITES", "type_objet": "methode", "capacites_codes": ["C5"], "methodes": ["M5"], "sources_inspiration": [], "status": "generated"}

\begin{fichemethode}{M5}{Étudier le sens de variation d'une suite}

\quandUtiliser{
Utiliser cette méthode lorsque l'énoncé demande de :
\begin{itemize}
  \item « montrer que la suite $(u_n)$ est croissante / décroissante » ;
  \item « étudier les variations de la suite $(u_n)$ » ;
  \item « montrer que $(u_n)$ est monotone à partir d'un certain rang ».
\end{itemize}

\medskip
\textbf{Arbre de décision — quelle méthode choisir ?}
\begin{itemize}
  \item Suite géométrique à termes positifs ou expression avec un quotient naturel $\rightarrow$ \textbf{Méthode B} (quotient).
  \item Suite définie par une expression de $n$ (polynomiale, rationnelle, etc.) $\rightarrow$ \textbf{Méthode A} (différence) ou \textbf{Méthode C} (fonction associée).
  \item Suite définie par récurrence $u_{n+1} = f(u_n)$ $\rightarrow$ \textbf{Méthode A} (différence).
\end{itemize}
}

\pasApas{
\textbf{Méthode A — Signe de $u_{n+1} - u_n$}
\begin{enumerate}
  \item Je calcule $u_{n+1} - u_n$ en développant et simplifiant.
  \item J'étudie le signe de cette expression pour tout $n \in \mathbb{N}$ (ou pour tout $n$ à partir d'un certain rang).
  \item Je conclus : si $u_{n+1} - u_n \geq 0$ pour tout $n$, alors $(u_n)$ est croissante ; si $u_{n+1} - u_n \leq 0$, alors $(u_n)$ est décroissante.
\end{enumerate}

\medskip
\textbf{Méthode B — Comparaison du quotient $\dfrac{u_{n+1}}{u_n}$ à $1$}
\begin{enumerate}
  \item Je vérifie que $u_n > 0$ pour tout $n \in \mathbb{N}$ (condition indispensable).
  \item Je calcule $\dfrac{u_{n+1}}{u_n}$ et je simplifie.
  \item Je compare ce quotient à $1$ pour tout $n \in \mathbb{N}$.
  \item Je conclus : si $\dfrac{u_{n+1}}{u_n} \geq 1$, alors $(u_n)$ est croissante ; si $\dfrac{u_{n+1}}{u_n} \leq 1$, alors $(u_n)$ est décroissante.
\end{enumerate}

\medskip
\textbf{Méthode C — Étude de la fonction associée}
\begin{enumerate}
  \item J'écris $u_n = f(n)$ où $f$ est une fonction définie sur $[0\,;\,+\infty[$.
  \item Je calcule la dérivée $f'$ et j'étudie son signe sur $[0\,;\,+\infty[$.
  \item Je conclus sur la monotonie de $(u_n)$ selon le signe de $f'$.
\end{enumerate}
}

\exempleRedige{
\textbf{Énoncé.} Soit $(u_n)$ la suite définie pour tout $n \in \mathbb{N}$ par $u_n = n^2 - 4n + 7$. Étudier les variations de $(u_n)$.

\medskip
On choisit la \textbf{méthode A} (expression polynomiale de $n$).

\medskip
\commentaireMarge{Calculer $u_{n+1} - u_n$ en substituant $n+1$ à $n$.}
Pour tout $n \in \mathbb{N}$ :
\[
  u_{n+1} = (n+1)^2 - 4(n+1) + 7 = n^2 + 2n + 1 - 4n - 4 + 7 = n^2 - 2n + 4.
\]

\commentaireMarge{Développer, puis simplifier pour faire apparaître une expression dont on peut étudier le signe.}
Donc :
\[
  u_{n+1} - u_n = (n^2 - 2n + 4) - (n^2 - 4n + 7) = 2n - 3.
\]

\commentaireMarge{Étudier le signe selon les valeurs de $n$. Ne pas oublier le quantificateur « pour tout $n$ ».}
On étudie le signe de $2n - 3$ :
\begin{itemize}
  \item Si $n \geq 2$, alors $2n \geq 4 > 3$, donc $2n - 3 > 0$, soit $u_{n+1} - u_n > 0$.
  \item Si $n = 0$, alors $u_1 - u_0 = -3 < 0$.
  \item Si $n = 1$, alors $u_2 - u_1 = -1 < 0$.
\end{itemize}

\commentaireMarge{Rédiger une conclusion précise, en distinguant les différents cas.}
\textbf{Conclusion.} La suite $(u_n)$ est \textbf{décroissante} pour $n \in \{0\,;\,1\}$ et \textbf{strictement croissante} pour tout entier $n \geq 2$. Elle n'est pas monotone sur $\mathbb{N}$ tout entier.

\medskip
\textit{Vérification numérique :} $u_0 = 7$, $u_1 = 4$, $u_2 = 3$, $u_3 = 4$, $u_4 = 7$. On constate bien $u_0 > u_1 > u_2$ puis $u_2 < u_3 < u_4$, ce qui est cohérent.
}

\pieges{
\textbf{Piège 1 — Oublier le quantificateur universel.}

\textit{Erreur :} « $u_1 - u_0 = 5 > 0$, donc la suite est croissante. »

\textit{Correction :} Vérifier $u_{n+1} - u_n \geq 0$ \textbf{pour tout} $n \in \mathbb{N}$. Un seul exemple ne constitue pas une preuve ; il faut un raisonnement général.

\medskip
\textbf{Piège 2 — Utiliser le quotient sans vérifier la positivité.}

\textit{Erreur :} Calculer $\dfrac{u_{n+1}}{u_n}$ alors que $u_n$ peut être nul ou négatif.

\textit{Correction :} La méthode B (quotient) n'est valide que si $u_n > 0$ \textbf{pour tout} $n \in \mathbb{N}$. Dans le cas contraire, utiliser la méthode A ou C.

\medskip
\textbf{Piège 3 — Conclure que la suite n'est pas monotone sans préciser les intervalles.}

\textit{Erreur :} « La suite n'est pas monotone. »

\textit{Correction :} Préciser les rangs : « La suite est décroissante pour $n \leq n_0$ et croissante pour $n \geq n_0$. »
}

\verifier{
\begin{itemize}
  \item Calculer les trois ou quatre premiers termes de la suite et vérifier qu'ils sont bien dans l'ordre annoncé (croissant ou décroissant).
  \item Si l'on a conclu à la croissance à partir du rang $n_0$, vérifier explicitement que $u_{n_0} \leq u_{n_0+1}$ par le calcul.
  \item Pour la méthode C : vérifier que $f$ est bien dérivable sur $[0\,;\,+\infty[$ avant de conclure.
\end{itemize}
}

\sEntrainer{\refExos{M5}}

\end{fichemethode}
